\section{Structure séquence}
\label{sec:extension-scripts-sequences}



\begin{figure}[htb!]
  \centering
  \includegraphics[height=40pt]{script-fil-rouge}
  \caption{Entrée du script pour le cas d'étude nº~1}
  \label{fig:reevaluation-scripts-sequence-exemple-cas-etude-1-entree}
\end{figure}
%
\begin{lstlisting}[style=script-style,
language=jerboa-script,
escapechar=\%,
caption=Script du cas d'étude nº~1,
label=lst:reevaluation-scripts-sequence-exemple-cas-etude-1
]
JerboaRuleResult resCube1 = @rule<extrusionFace>(face#0);
// resCube1#4#7 : (colonne 4; rangee 7)
JerboaRuleResult resVertex = @rule<insertionSommet>(resCube1#4#7);
JerboaRuleResult resCube2 = @rule<extrusionFaceVolume>(resCube1#5#7);
JerboaRuleResult resTri = @rule<triangulationFace>(resCube1#1#5);
JerboaRuleResult resContraction = @rule<contractionFace>(resCube2#6#1);
JerboaRuleResult resChanfrein = @rule<chanfreinSommet>(resCollapse#0#1);
@rule<coloration>(resTri#0#2,blue);
@rule<coloration>(resChamfer#0#0,red);
return null;
\end{lstlisting}

Dans cette section, nous proposons les premières modifications nécessaires aux
spécifications paramétriques ainsi qu'aux DAGs afin de supporter la réévaluation
des scripts. %
Prenons pour exemple le script
(liste~\ref{lst:reevaluation-scripts-sequence-exemple-cas-etude-1}) qui prend
une face en entrée
(figure~\ref{fig:reevaluation-scripts-sequence-exemple-cas-etude-1-entree}) et
qui produit le cas d'étude nº~1 du
chapitre~\ref{cha:reevaluation-base-regles}. %



\subsection{Spécification paramétrique}
\label{sec:reevaluation-scripts-sequence-spec-param}



La représentation en boîte fermée présente plusieurs inconvénients. %
En particulier, l'impossibilité d'accéder aux règles utilisées et à leurs
paramètres topologiques. %
En effet, ne pas pouvoir accéder aux règles, implique de ne pas pouvoir procéder
à leurs analyses. %
Les développeurs d'opérations doivent alors enregistrer manuellement les
évènements topologiques. %
De plus, ne pas avoir accès aux paramètres topologiques signifie proposer un
suivi dégradé de l'évolution des orbites et donc un appariement imprévisible des
orbites lors de la réévaluation. %
Or, cette démarche est incompatible avec celle que nous proposons dans nos
travaux. %
Nous proposons donc d'enregistrer les opérations de modélisation de manière
transparente avec le modèle de boîte ouverte. %

\begin{lstlisting}[style=script-style,
escapechar=\%,
caption=Spécification paramétrique du cas d'étude nº~1 avec un script,
label=lst:reevaluation-scripts-sequence-spec-param-cas-etude-1]
1-creationCarre()
2-casEtude1(pn1)
  Sequence:
    1-extrusionFace(pn2)
    2-insertionSommet(pn3)
    3-extrusionFaceVolume(pn4)
    4-triangulationFace(pn5)
    5-contractionFace(pn6)
    6-chanfreinSommet(pn7)
    7-coloration(pn8)
    8-coloration(pn9)
\end{lstlisting}

Avec la représentation boîte ouverte, nous avons accès aux instructions des
scripts. %
Il nous est alors possible d'enregistrer dans les spécifications paramétriques
le contenu dont nous avions besoin jusque-là pour les règles, à savoir leurs
noms et leurs paramètres topologiques. %
Toutefois, dans notre approche, nous distinguons le contenu d'un script des
autres opérations enregistrées dans une spécification paramétrique. %
L'enregistrement d'un script ouvre un bloc, avec indentation, où sont nommées
les structures et leur contenu. %
Par exemple, dans la spécification paramétrique représentée dans la
liste~\ref{lst:reevaluation-scripts-sequence-spec-param-cas-etude-1}, un bloc
\texttt{Sequence} est enregistré sous l'enregistrement du script
\texttt{casEtude1} et contient l'ensemble des opérations appelées. %
\todo{Reprendre ici (tout le paragraphe, ça part dans tous les sens)}

%% ce que l'on peut faire
%%% connaître les règles et les paramètres topologiques pour l'appariement
%%% prévenir une erreur d'application grâce au suivi des orbites
%% contraintes
%%% les mêmes que les règles + la construction manuelle de la sortie (jerboaruleresult)

% Les boîtes ouvertes le mode de représentation que nous préférons, car celles-ci
% permettent au système de réévaluation d'avoir connaissances des règles (et
% scripts) appelés en interne d'un script et de faire l'appariement de leurs
% paramètres topologiques respectifs. %
% Cependant, cette deuxième approche requiert l'extension des structures de
% données nous permettant de représenter les diverses étapes de modélisation, à
% savoir la spécification paramétrique ainsi que les DAGs d'évaluation et de
% réévaluation. %
% Dans cette configuration, il est possible de gérer les interruptions de
% réévaluation puisque le système connaît les règles d'un script et leurs
% paramètres topologiques. %
% Pour les développeurs, la contrainte est similaire aux règles Jerboa. %
% En effet, l'application d'une règle Jerboa retourne un résultat d'application
% sous la forme d'un tableau de brins instanciés à partir du membre droit. %
% Pour les scripts, les développeurs doivent aussi construire cette sortie. %

\subsection{DAGs}
\label{sec:reevaluation-scripts-sequence-extension-dags}

\todo{DAG évaluation avec bloc boîte ouverte}

L'utilisation d'une boîte ouverte pour représenter les opérations internes à un
script nous permet de reconstituer les historiques d'évolutions de chacun des
paramètres topologiques dans des DAGs. %
Afin de faire la distinction entre les opérations hors scripts et les opérations
dans les scripts, le modèle de boîte ouverte doit être reportée dans les DAGs
notamment parce que, tout comme les règles, les scripts peuvent référencer une
entité topologique comme paramètre. %
Ce paramètre conditionne alors l'exécution du script et l'ouverture de la
boîte. %


%%% Local Variables:
%%% mode: LaTeX
%%% TeX-master: "../../../main"
%%% End:
